\documentclass[11pt]{article}
\usepackage[margin=1in]{geometry}
\usepackage{hyperref}
\usepackage{enumitem}
\usepackage{titlesec}
\usepackage{fancyhdr}
\usepackage{booktabs}

% Header/Footer
\pagestyle{fancy}
\fancyhf{}
\rhead{CS 4300/5300 Software Engineering}
\lhead{Sprint 0-0}
\rfoot{Page \thepage}

% Section formatting
\titleformat{\section}{\Large\bfseries}{\thesection}{1em}{}
\titleformat{\subsection}{\large\bfseries}{\thesubsection}{1em}{}

\begin{document}

\begin{center}
{\LARGE\textbf{CS 4300/5300 Software Engineering}}\\[0.5em]
{\Large Sprint 0-0: Project Idea}\\[1em]
\end{center}

\noindent\textit{Sprint 0 is preparation work. It runs before your first real iteration, so the Common Sprint Requirements document does not apply yet---there is no demo, no sprint tag, and no peer evaluation for Sprint 0-0.}

\section{Purpose}

You cannot build software until you know who you are building it for. Sprint 0-0 is where your team decides what it wants to build and describes it clearly enough that someone else can tell whether it is worth building.

Your whole team submits \textbf{one} project idea, through the Google Form linked below. One submission per team.

\section{Constraints}

Whatever you propose must satisfy both of these:
\begin{itemize}
    \item it uses a \textbf{database backend}; and
    \item it integrates at least one \textbf{third-party API}.
\end{itemize}

You will be building this across four sprints, and another team will act as your customer. Pick a problem that someone other than your team actually has.

\section{Sprint 0-0 Objectives}

\begin{center}
\begin{tabular}{lc}
\toprule
\textbf{Requirement} & \textbf{Points} \\
\midrule
Project title & 5 pts \\
Project description (at least 100 words) & 25 pts \\
Objectives and goals (at least 5) & 20 pts \\
Target audience & 15 pts \\
Key features (at least 10) & 20 pts \\
Available APIs (3 that are free to use) & 15 pts \\
\midrule
\textbf{Total} & \textbf{100 pts} \\
\bottomrule
\end{tabular}
\end{center}

\section{Detailed Requirements}

\subsection{Project Title (5 pts)}
A working title for the application. It can change later.

\subsection{Project Description (25 pts)}
At least \textbf{100 words} describing what the application does.

Describe the \textbf{problem}, not just the feature list. ``Users can log in and see a dashboard'' is a feature. ``Students lose track of which assignments are actually urgent'' is a problem. Cover the core problem the application solves, who has that problem, and why solving it matters to them.

\subsection{Objectives and Goals (20 pts)}
At least \textbf{five} specific goals or outcomes the project seeks to achieve.

State outcomes, not activities. ``Write the login page'' is a task. ``A user can recover their account without contacting support'' is an outcome.

\subsection{Target Audience (15 pts)}
Who this is for, described as a \textbf{role} rather than a person. Name your primary users and say how they differ from anyone else who touches the system.

\subsection{Key Features (20 pts)}
At least \textbf{ten} features. One line each is fine. These are the raw material for the user stories you will write in Sprint 0-2, so favor things a user can observably do over internal implementation details.

\subsection{Available APIs (15 pts)}
Find \textbf{three} third-party APIs that are free to use and could serve your project. For each, name it and state in one line what your application would get from it.

Confirm that each one is actually free at the usage level you need---some APIs are free to sign up for and not free to call. You are not committing to all three; you are demonstrating that your project has a viable integration path.

\section{Submission Requirements}

\subsection*{Google Form (1 submission per team)}
Submit through the course project idea form:

\begin{center}
\url{https://docs.google.com/forms/d/e/1FAIpQLSelTBtrf8bdzrYCIhtc_LasULG46gb5awVTDBmu3nWoWmYMzQ/viewform}
\end{center}

The form asks for your email and your team number in addition to the content above. Decide as a team who submits, and make sure the rest of the team has seen the text before it is sent---you cannot edit it afterward.

\subsection*{AI Documentation}
If you used AI while preparing this submission, document where and when, and include the transcript URL for each instance.

\end{document}
